% MRV -- Model Risk Validation Report
% All text lives here. report.py only substitutes data values.
%
% Placeholders: {{KEY}} — replaced with data from JSON
% Conditionals: %% IF_xxx ... %% ELSE ... %% ENDIF — selected by report.py
% Asset block:  %% BEGIN_ASSET ... %% END_ASSET — repeated per asset

\documentclass[11pt,a4paper]{article}

\usepackage[utf8]{inputenc}
\usepackage[T1]{fontenc}
\usepackage[margin=2.5cm]{geometry}
\usepackage{booktabs}
\usepackage{graphicx}
\usepackage{xcolor}
\usepackage{fancyhdr}
\usepackage{titlesec}
\usepackage{hyperref}
\usepackage{float}
\usepackage{caption}
\usepackage{enumitem}
\usepackage{amsmath}
\usepackage{colortbl}
\usepackage{tabularx}
\usepackage{tikz}

\definecolor{mrvblue}{RGB}{0,51,102}
\definecolor{mrvgray}{RGB}{100,100,100}
\definecolor{mrvlightgray}{RGB}{240,240,240}
\definecolor{mrvred}{RGB}{180,30,30}
\definecolor{mrvgreen}{RGB}{30,120,50}
\definecolor{mrvredbg}{RGB}{253,235,235}
\definecolor{mrvgreenbg}{RGB}{232,248,237}

\newcommand{\statuspass}{\tikz\draw[fill=mrvgreen,draw=none] (0,0) circle (0.35em); \textcolor{mrvgreen}{\textbf{PASS}}}
\newcommand{\statusfail}{\tikz\draw[fill=mrvred,draw=none] (0,0) circle (0.35em); \textcolor{mrvred}{\textbf{FAIL}}}

\pagestyle{fancy}
\fancyhf{}
\fancyhead[L]{\small\color{mrvgray}MRV --- Model Risk Validation Report}
\fancyhead[R]{\small\color{mrvgray}{{DATE}}}
\fancyfoot[L]{\footnotesize\color{mrvgray}CONFIDENTIAL \enspace$\cdot$\enspace \href{https://github.com/modelguard-lab/mrv-lib}{mrv-lib}}
\fancyfoot[R]{\small\color{mrvgray}\thepage}
\renewcommand{\headrulewidth}{0.4pt}
\renewcommand{\footrulewidth}{0.2pt}

\titleformat{\section}{\large\bfseries\color{mrvblue}}{\thesection}{1em}{}
\titleformat{\subsection}{\normalsize\bfseries\color{mrvblue}}{\thesubsection}{1em}{}
\hypersetup{colorlinks=true, linkcolor=mrvblue, urlcolor=mrvblue}

\begin{document}

% =====================================================================
% COVER
% =====================================================================
\begin{titlepage}
\centering
\vspace*{2cm}

{\Huge\bfseries\color{mrvblue} Model Risk Validation Report}\\[0.6cm]
{\LARGE Representation Invariance Assessment}\\[1.5cm]

%% IF_PARTITION_PASS
\fcolorbox{mrvgreen}{mrvgreenbg}{%
  \parbox{0.7\textwidth}{\centering
    \vspace{0.4cm}
    {\Large\textbf{Partition Stability:} \statuspass}\\[0.2cm]
    {\large Mean ARI = {{OVERALL_MEAN_ARI}} \quad ($\geq$ {{ARI_THRESHOLD}})}
    \vspace{0.4cm}
  }%
}
%% ELSE
\fcolorbox{mrvred}{mrvredbg}{%
  \parbox{0.7\textwidth}{\centering
    \vspace{0.4cm}
    {\Large\textbf{Partition Stability:} \statusfail}\\[0.2cm]
    {\large Mean ARI = {{OVERALL_MEAN_ARI}} \quad ($<$ {{ARI_THRESHOLD}})}
%% IF_ORDERING_PASS
    \\[0.4cm]
    {\Large\textbf{Ordering Stability:} \statuspass}\\[0.2cm]
    {\large Mean Spearman = {{OVERALL_MEAN_SPEARMAN}} \quad ($\geq$ {{SPEARMAN_THRESHOLD}})}
%% ENDIF
    \vspace{0.4cm}
  }%
}
%% ENDIF

\vspace{1.5cm}

\begin{tabular}{ll}
\toprule
\textbf{Date}       & {{DATE}} \\
\textbf{Scope}      & {{ASSET_LIST}} --- representation stability \\
\textbf{Model}      & {{MODEL}} ($K = {{N_STATES}}$) \\
\textbf{Period}     & {{DATE_RANGE}} \\
\textbf{Framework}  & SR 11-7 / Basel IV FRTB \\
\bottomrule
\end{tabular}

\vspace{2cm}

\begin{tabular}{p{4.5cm}p{4.5cm}p{4.5cm}}
\toprule
\textbf{Prepared by} & \textbf{Reviewed by} & \textbf{Approved by} \\
\midrule
& & \\[1.2cm]
\dotfill & \dotfill & \dotfill \\
Name & Name & Name \\
Date: \dotfill & Date: \dotfill & Date: \dotfill \\
\bottomrule
\end{tabular}

\vfill
{\footnotesize\color{mrvgray} CONFIDENTIAL --- For internal risk management use only.}
\end{titlepage}


\tableofcontents
\newpage


% =====================================================================
% 1. DASHBOARD
% =====================================================================
\section{Dashboard}

\noindent
At-a-glance results for all assets under the representation invariance check.
A \statuspass{} requires mean cross-representation ARI at or above the 0.65
threshold established by Steinley (2004) as the minimum for acceptable
cluster recovery quality.

\vspace{0.5cm}
{{DASHBOARD_TABLE}}

\vspace{0.3cm}
\noindent
%% IF_PARTITION_FAIL_ORDERING_PASS
\textbf{Key finding:} Categorical regime labels are \textbf{not stable}
across feature representations (ARI below {{ARI_THRESHOLD}}). However, coarse
\textbf{risk ordering is preserved} (Spearman above {{SPEARMAN_THRESHOLD}}), meaning
ordinal risk signals remain reliable even when exact state boundaries shift.
%% ELIF_PARTITION_FAIL
\textbf{Key finding:} Both partition labels and risk ordering are unstable
across representations. Regime model outputs should not be used for
risk decisions without further investigation.
%% ELSE
\textbf{Key finding:} Regime labels show acceptable stability
across feature representations at both partition and ordering levels.
%% ENDIF


% =====================================================================
% 2. EXECUTIVE SUMMARY
% =====================================================================
\section{Executive Summary}

\subsection{What was tested}
We tested whether the regime labels produced by a {{MODEL}} ($K={{N_STATES}}$)
model remain consistent when the input features are changed.
The following economically reasonable feature sets were compared:

{{FACTOR_SETS}}

\subsection{What we found}
\begin{enumerate}[leftmargin=2em]
%% IF_PARTITION_FAIL
\item \textbf{Regime labels are unstable.} The same observation window can be
      classified differently under different feature sets. Mean ARI across
      all assets is {{OVERALL_MEAN_ARI}}, below the {{ARI_THRESHOLD}} threshold
      (Steinley 2004).
%% ELSE
\item \textbf{Partition stability: PASS.} Mean ARI = {{OVERALL_MEAN_ARI}}
      ($\geq$ {{ARI_THRESHOLD}}).
%% ENDIF
%% IF_ORDERING_PASS
\item \textbf{Risk ordering is stable.} Despite label instability, the
      ranking of states from lowest to highest risk is consistent across
      representations (mean Spearman = {{OVERALL_MEAN_SPEARMAN}},
      $\geq$ {{SPEARMAN_THRESHOLD}}).
      The ordinal signal (``which state is riskiest?'') survives; the
      categorical signal (``which state are we in?'') does not.
%% ELSE
\item \textbf{Risk ordering is also unstable.} Mean Spearman =
      {{OVERALL_MEAN_SPEARMAN}} ($<$ {{SPEARMAN_THRESHOLD}}).
%% ENDIF
\end{enumerate}

\subsection{What should be done}
\begin{table}[H]
\centering
\renewcommand{\arraystretch}{1.3}
\begin{tabularx}{\textwidth}{c X l}
\toprule
\textbf{\#} & \textbf{Action} & \textbf{Priority} \\
\midrule
1 & Run regime inference with $\geq$3 feature sets; flag disagreement. & High \\
2 & Migrate from categorical labels to ordinal risk scores. & High \\
3 & Add representation sensitivity to model risk disclosures (SR 11-7). & Medium \\
4 & Schedule quarterly re-validation. & Medium \\
\bottomrule
\end{tabularx}
\caption{Recommended actions.}
\end{table}


% =====================================================================
% 3. METHODOLOGY
% =====================================================================
\section{Methodology}

A Gaussian Mixture Model with $K={{N_STATES}}$ states is fitted independently on
each feature set via Expectation--Maximization (seed = 42). Labels
are assigned by maximum posterior probability. Agreement between
label sequences is measured by the Adjusted Rand Index (ARI,
Hubert \& Arabie 1985), which is permutation-invariant and corrected
for chance.

\medskip\noindent
Following Steinley (2004), ARI $\geq$ 0.90 indicates excellent recovery,
ARI $\geq$ 0.65 indicates acceptable recovery, and
ARI $<$ 0.65 signals material representation dependence.

\medskip\noindent
All features are computed from log returns and normalized via rolling
$z$-score (120-bar lookback) to ensure comparability across representations.


% =====================================================================
% 4. RESULTS BY ASSET
% =====================================================================
\section{Results by Asset}

%% BEGIN_ASSET
\subsection{{{ASSET_NAME}}}

\begin{minipage}[t]{0.48\textwidth}
\begin{table}[H]
\centering
\caption{Pairwise ARI --- {{ASSET_NAME}}}
{{ARI_TABLE}}
\end{table}

\vspace{0.3cm}
\begin{table}[H]
\centering
\caption{Summary --- {{ASSET_NAME}}}
\begin{tabular}{ll}
\toprule
Observations  & {{N_OBS}} \\
Factor sets   & {{N_FACTOR_SETS}} \\
Mean ARI      & \textcolor{{{ARI_COLOR}}}{\textbf{{{MEAN_ARI}}}} \\
Min ARI       & {{MIN_ARI}} \\
Spearman      & {{MEAN_SPEARMAN}} \\
\bottomrule
\end{tabular}
\end{table}
\end{minipage}
\hfill
\begin{minipage}[t]{0.48\textwidth}
\vspace{0pt}
\begin{figure}[H]
\centering
\includegraphics[width=\textwidth]{{{HEATMAP_FILE}}}
\caption{ARI heatmap --- {{ASSET_NAME}}}
\end{figure}
\end{minipage}

\vspace{0.3cm}
\noindent
{{FINDING_TEXT}}

%% END_ASSET


% =====================================================================
% 5. CONCLUSION
% =====================================================================
\newpage
\section{Conclusion \& Recommendations}

\subsection{Risk statement}

%% IF_PARTITION_FAIL_ORDERING_PASS
The regime model under review ({{MODEL}}, $K={{N_STATES}}$) produces
\textbf{partition labels that are materially representation-dependent}
(mean ARI = {{OVERALL_MEAN_ARI}}, below {{ARI_THRESHOLD}}).
However, \textbf{coarse risk ordering is stable}
(mean Spearman = {{OVERALL_MEAN_SPEARMAN}}, above {{SPEARMAN_THRESHOLD}}).
Risk managers should not rely on categorical state labels for trading
decisions, but ordinal risk scores (``risk is elevated'') remain defensible.
%% ELIF_PARTITION_FAIL
The regime model ({{MODEL}}, $K={{N_STATES}}$) produces labels that are
\textbf{materially dependent on feature representation} at both
partition and ordering levels. Regime-based signals should be suspended
pending further investigation.
%% ELSE
The regime model ({{MODEL}}, $K={{N_STATES}}$) produces labels with acceptable
cross-representation stability (mean ARI = {{OVERALL_MEAN_ARI}}).
%% ENDIF

As required by the Federal Reserve's SR 11-7 guidance on model risk
management, this representation sensitivity constitutes a material model
risk that should be disclosed and mitigated.

\subsection{Recommended actions}

\begin{table}[H]
\centering
\renewcommand{\arraystretch}{1.3}
\begin{tabularx}{\textwidth}{c X l}
\toprule
\textbf{\#} & \textbf{Action} & \textbf{Priority} \\
\midrule
1 & Run regime inference with $\geq$3 feature sets; flag disagreement. & High \\
2 & Migrate from categorical labels to ordinal risk scores. & High \\
3 & Add representation sensitivity to model risk disclosures (SR 11-7). & Medium \\
4 & Schedule quarterly re-validation. & Medium \\
\bottomrule
\end{tabularx}
\caption{Remediation plan.}
\end{table}


% =====================================================================
% APPENDIX
% =====================================================================
\newpage
\appendix

\section{Data Lineage}

\begin{table}[H]
\centering
\begin{tabularx}{\textwidth}{l X}
\toprule
\textbf{Item} & \textbf{Detail} \\
\midrule
Data source     & Interactive Brokers (5-minute OHLCV bars) \\
Returns         & Log returns: $r_t = \log(P_t / P_{t-1})$ \\
Normalization   & Rolling $z$-score, 120-bar lookback \\
Missing values  & Excluded mechanically (warm-up period) \\
\bottomrule
\end{tabularx}
\end{table}

\section{Software \& Reproducibility}

\begin{table}[H]
\centering
\begin{tabularx}{\textwidth}{l X}
\toprule
\textbf{Item} & \textbf{Detail} \\
\midrule
Software        & mrv-lib (Model Risk Validator) \\
Random seed     & 42 \\
ARI threshold   & {{ARI_THRESHOLD}} (Steinley 2004) \\
Config file     & \texttt{config.yaml} (attached) \\
\bottomrule
\end{tabularx}
\end{table}

\end{document}
